\mxunit{premiere-c02}{Second degré : formes et choix de méthode}\label{C02}
\begin{keypoint}[Le parcours du chapitre]
\textbf{Objectifs.} Choisir une forme ; résoudre une équation ou une inéquation ; lire un extremum ; démontrer la formule des racines.

\textbf{Prérequis.} Identités remarquables, produit nul, signe d’un produit, racines carrées.

\textbf{Organisation.} Préparer les automatismes, chercher, formaliser, puis réinvestir.\end{keypoint}
\section*{Diagnostic : \mxcount{diag} questions, 10 minutes}\label{C02-D}
Répondre sans consulter le cours. Indiquer les questions pour lesquelles une aide est nécessaire. Ce diagnostic ne donne pas lieu à une note.
\begin{enumerate}
\item Développer $(x-3)^2$.
\item Factoriser $x^2-16$.
\item Résoudre $(x-2)(x+5)=0$.
\item Résoudre $x^2=9$ dans $\R$.
\item Donner le signe de $2x-6$.
\item Déterminer le minimum de $(x-1)^2+4$.
\mxcounted{diag}
\end{enumerate}
\paragraph{Après le diagnostic.} Reprendre les notions signalées avec le professeur ; refaire une question analogue, puis expliquer la règle utilisée.
\paragraph{Oral et logique.} Présenter deux méthodes pour factoriser un trinôme. Dans la démonstration par discriminant, nommer les trois cas et expliquer pourquoi ils sont exhaustifs.
\mxnewpage
\section*{Chercher avant de formaliser}\label{C02-RECH}
Un rectangle a un périmètre de 24 unités. Choisir plusieurs largeurs strictement positives donnant un rectangle. Pour les largeurs 2, 4, 6, 8 et 10, calculer la longueur et l’aire. Quelle largeur semble donner la plus grande aire ? Formuler une conjecture, puis chercher une justification valable pour toutes les largeurs admissibles.
\paragraph{Organisation.} Recherche individuelle 5 min, confrontation à deux 5 min, mise en commun 5 min. Garder une trace d’un essai, même non abouti. Une conjecture doit ensuite être justifiée.
\mxnewpage
\section{Trois formes, trois lectures}\label{C02-U1}
\begin{keypoint}[Définition et résultats]
Une fonction du second degré s’écrit $f(x)=ax^2+bx+c$ avec $a\ne0$. La forme développée facilite le calcul de $f(0)$ ; la forme canonique $a(x-\alpha)^2+\beta$ révèle sommet et extremum ; une forme factorisée $a(x-r_1)(x-r_2)$ révèle racines et signe. Les racines sont les solutions de $f(x)=0$.
\end{keypoint}
\paragraph{Statut et justification.} La complétion du carré repose sur $x^2+2ux=(x+u)^2-u^2$. L’existence de la forme canonique est démontrée par développement ; son calcul général n’est pas un automatisme exigible.
\begin{example}$f(x)=x^2-6x+5=(x-3)^2-4=(x-1)(x-5)$. Le minimum est $-4$ en 3 ; les racines sont 1 et 5 ; l’ordonnée à l’origine vaut 5.\end{example}\label{C02-EX1}
\mxinclude{figures/C02-U1}
\begin{pitfall}Une expression factorisée ne s’obtient pas en mettant séparément chacun des termes entre parenthèses.\end{pitfall}
\begin{application}Donner trois formes de $g(x)=x^2+2x-3$. Lire ses racines et son minimum.\end{application}\label{C02-AP1}
\section{Discriminant et résolution}\label{C02-U2}
\begin{keypoint}[Définition et résultats]
Pour $ax^2+bx+c=0$ avec $a\ne0$, on pose $\Delta=b^2-4ac$. Si $\Delta<0$, aucune racine réelle ; si $\Delta=0$, une racine $-\dfrac{b}{2a}$ ; si $\Delta>0$, deux racines $\dfrac{-b-\sqrt\Delta}{2a}$ et $\dfrac{-b+\sqrt\Delta}{2a}$. Pour ordonner les racines, tenir compte du signe de a.
\end{keypoint}
\paragraph{Statut et justification.} En multipliant l’équation par $4a\ne0$, on obtient $(2ax+b)^2=\Delta$. Un carré est positif ou nul : les trois cas suivent. Si $\Delta>0$, $2ax+b=\pm\sqrt\Delta$. La disjonction des cas évite de prendre la racine carrée d’un nombre négatif.
\begin{example}$2x^2+x-3=0$ : $\Delta=1+24=25$ ; $x=\dfrac{-1\pm5}{4}$, soit $1$ et $-\dfrac{3}{2}$. Vérification : $(x-1)(2x+3)=2x^2+x-3$.\end{example}\label{C02-EX2}
\begin{pitfall}Le dénominateur est $2a$, pas 2 ; le signe de b est pris avec ses parenthèses dans $b^2$.\end{pitfall}
\begin{application}Résoudre $3x^2-6x+3=0$ et $x^2+2x+5=0$.\end{application}\label{C02-AP2}
\section{Factoriser sans automatiser le discriminant}\label{C02-U3}
\begin{keypoint}[Définition et résultats]
Si deux racines réelles existent (éventuellement confondues, alors comptées deux fois), leur somme vaut $-\dfrac{b}{a}$ et leur produit $\dfrac{c}{a}$. Une racine apparente ou une identité remarquable permet souvent une factorisation plus courte. Les polynômes de degré 2 s’annulant en deux réels distincts r et s sont exactement $a(x-r)(x-s)$ avec $a\ne0$.
\end{keypoint}
\paragraph{Statut et justification.} En développant $a(x-r)(x-s)=ax^2-a(r+s)x+ars$, l’identification des coefficients donne les relations. Réciproquement, si r et s ont cette somme et ce produit, le développement retrouve le polynôme.
\begin{example}Posons $P(x)=x^2-7x+12$. Les nombres 3 et 4 ont pour somme 7 et pour produit 12, donc $P(x)=(x-3)(x-4)$. Autre question : déterminer une fonction $f$ de racines 3 et 4 avec $f(0)=24$. Écrivons $f(x)=a(x-3)(x-4)$. Alors $12a=24$, donc $a=2$ : $f$ est différente de $P$.\end{example}\label{C02-EX3}
\begin{pitfall}La somme des racines vaut $-\dfrac{b}{a}$ ; ne pas oublier le signe moins. Les racines ne déterminent pas a : il faut une condition supplémentaire, par exemple la valeur de f en un point.\end{pitfall}
\begin{application}Déterminer le polynôme f de degré 2 tel que $f(1)=f(4)=0$ et $f(0)=-8$.\end{application}\label{C02-AP3}
\section{Signe, inéquations et optimisation}\label{C02-U4}
\begin{keypoint}[Définition et résultats]
Pour deux racines distinctes $r<s$, le trinôme a le signe de a en dehors de $\intff{r}{s}$ et le signe opposé entre les racines. Il est nul aux racines. Sans racine réelle, il garde le signe de a. Avec une racine double, il garde le signe de a hors de celle-ci. L’optimisation se fait sur le domaine du problème, qui peut limiter l’accès au sommet.
\end{keypoint}
\paragraph{Statut et justification.} Le tableau de signes des facteurs $a(x-r)(x-s)$ donne le résultat. La forme canonique compare $f(x)$ à $\beta$ grâce au signe de $a(x-\alpha)^2$.
\begin{example}$-2(x-1)(x-5)\geq0$ équivaut à $x\in\intff{1}{5}$. Pour $A(x)=x(12-x)$ sur $\intff{0}{12}$, $A(x)=36-(x-6)^2\leq36$, avec égalité en 6. Un rectangle de périmètre 24 m a donc une aire maximale de 36 m$^2$ lorsqu’il est carré.\end{example}\label{C02-EX4}
\begin{pitfall}Dans une inéquation stricte, les racines sont exclues. Vérifier les bornes géométriques avant de conclure.\end{pitfall}
\begin{application}Résoudre $x^2-4x-5<0$ puis maximiser $x(8-x)$ sur $\intff{0}{3}$.\end{application}\label{C02-AP4}
\mxnewpage
\section*{Parcours d’exercices}\label{C02-PARC}
Commencer par les exercices 1 et 2. Choisir ensuite un exercice 3 ou 4 sans aide. Les exercices 5 à 7 sont une réserve de transfert, de recherche ou de programmation : le professeur en désigne un. L’exercice 8 se reprend à distance.
\begingroup
\mxexercisehead{C02-E1}{Compléter un carré}
On considère la fonction $f$ définie sur $\R$ par $f(x)=x^2-8x+7$.
\begin{enumerate}[label=\arabic*.]
\item Compléter le carré pour écrire $f(x)$ sous la forme $(x-4)^2+k$, où $k$ est un réel à déterminer.
\item En déduire une factorisation de $f(x)$ à l’aide d’une identité remarquable.
\item Dresser le tableau de signes de $f$, puis résoudre dans $\R$ l’inéquation $f(x)\leq0$.
\end{enumerate}
\mxexercisehead{C02-E2}{Discuter sans calcul lourd}
On étudie les fonctions $f(x)=x^2+4$, $g(x)=(x+2)^2$ et $h(x)=-3(x-1)(x+4)$, toutes définies sur $\R$.
\begin{enumerate}[label=\arabic*.]
\item Pour chaque fonction, déterminer les éventuelles racines réelles en exploitant directement la forme donnée.
\item Préciser le signe de chaque fonction sur $\R$. Indiquer explicitement les valeurs où elle s’annule.
\item Expliquer pourquoi le calcul d’un discriminant n’est nécessaire dans aucun des trois cas.
\end{enumerate}
\mxexercisehead{C02-E3}{Choisir une stratégie}
On souhaite résoudre dans $\R$ les trois équations suivantes :\[x^2-9=0,\qquad x^2-9x+20=0,\qquad 2x^2-3x-1=0.\]
\begin{enumerate}[label=\arabic*.]
\item Pour chaque équation, choisir une méthode : identité remarquable, recherche de deux nombres de somme et produit connus, ou discriminant. Expliquer ce choix.
\item Résoudre les trois équations et donner les ensembles de solutions sous forme exacte.
\item Contrôler les solutions des deux premières équations par substitution. Pour la troisième, contrôler la somme et le produit des racines obtenues.
\end{enumerate}
\mxexercisehead{C02-E4}{Retrouver une parabole}
La courbe d’une fonction polynôme $f$ de degré 2 coupe l’axe des abscisses aux points d’abscisses $-2$ et 3. Elle passe également par le point $A\pt{1;12}$.
\begin{enumerate}[label=\arabic*.]
\item Déterminer une expression de $f(x)$ à partir de ces trois informations. Préciser la forme choisie et pourquoi elle convient.
\item Développer l’expression trouvée. En déduire l’ordonnée à l’origine de la courbe.
\item Vérifier que la fonction obtenue satisfait les trois informations initiales.
\end{enumerate}
\mxexercisehead{C02-E5}{Un potager contre un mur}
Un jardinier dispose de 20 m de grillage pour entourer trois côtés d’un potager rectangulaire. Le quatrième côté longe un mur et ne nécessite pas de grillage. On note $x$ la longueur, en mètres, de chacun des deux côtés perpendiculaires au mur. Toute la longueur de grillage est utilisée.
\begin{enumerate}[label=\arabic*.]
\item Faire un schéma et exprimer la longueur du côté parallèle au mur en fonction de $x$. Préciser les valeurs de $x$ pour lesquelles le rectangle n’est pas aplati.
\item Exprimer l’aire $A(x)$, en mètres carrés. Déterminer sa valeur maximale et les dimensions correspondantes, en choisissant la méthode de votre choix.
\item Contrôler que ces dimensions utilisent exactement 20 m de grillage.
\end{enumerate}
\mxinclude{figures/C02-E5}
\mxexercisehead{C02-E6}{Un paramètre}
Pour chaque réel $m$, on considère l’équation d’inconnue réelle $x$ :\[x^2-2x+m=0.\]Le nombre $m$ est un paramètre : il reste fixé lorsque l’on résout l’équation.
\begin{enumerate}[label=\arabic*.]
\item Déterminer les valeurs de $m$ pour lesquelles l’équation admet deux solutions réelles distinctes.
\item Déterminer les valeurs de $m$ pour lesquelles elle admet une unique solution réelle, puis celles pour lesquelles elle n’en admet aucune.
\item Donner la solution dans le cas où elle est unique. Expliquer pourquoi les trois cas envisagés couvrent toutes les valeurs de $m$.
\end{enumerate}
\mxexercisehead{C02-E7}{Lien avec les suites}
On définit une suite par $u_n=n^2-5n+4$ pour tout entier naturel $n$. On cherche tous les rangs auxquels le terme de la suite est strictement négatif.
\begin{enumerate}[label=\arabic*.]
\item Résoudre d’abord dans $\R$ l’inéquation $x^2-5x+4<0$.
\item En déduire tous les entiers naturels $n$ tels que $u_n<0$. Justifier pourquoi un intervalle réel ne constitue pas ici la réponse finale.
\end{enumerate}
\mxexercisehead{C02-E8}{Reprise différée}
Une élève étudie la fonction $f$ définie sur $\R$ par $f(x)=2(x+3)^2-5$. Elle écrit : « Le minimum vaut $-5$ et il est atteint pour $x=3$. »
\begin{enumerate}[label=\arabic*.]
\item Indiquer ce qui est correct et ce qui est faux dans cette phrase. La réécrire correctement en justifiant votre réponse.
\item Résoudre dans $\R$ l’inéquation $f(x)\leq-3$, puis contrôler les deux bornes de l’ensemble obtenu.
\end{enumerate}
\paragraph{Travail personnel.} Effectuer uniquement la tâche désignée à l’issue de la séance, dans la progression datée : 10–20 min maximum. Les exercices dépendent des notions déjà étudiées ; certaines questions sont donc différées. Reprendre l’exercice 8 une semaine après le chapitre (10 min). Si le temps prévu est dépassé, noter la difficulté et apporter les essais ; ne pas accumuler toute la réserve.
\endgroup
\mxneed{12}\section*{À faire dans le manuel CQFD}
Ces exercices se traitent dans votre manuel. La progression indique la séance ; ils remplacent le travail personnel prévu ce jour-là (10 à 20 minutes au total).
\paragraph{CQFD, exercice 43 p. 99, questions a et c.}\label{C02-M1} Trouver d’abord une racine évidente, puis l’autre à l’aide du produit des racines.
\paragraph{CQFD, exercice 49 p. 99, questions a et b.}\label{C02-M2} Traduire la somme et le produit par une équation ; ne pas supposer d’avance que les nombres existent.
\begingroup\small
\mxnewpage\section*{QCM 1 : comprendre et choisir}\label{C02-Q1}
\mxCount{qcm1} questions pédagogiques à réponse unique. Choisir une réponse par question ; expliquer ensuite un distracteur. Ce format de chapitre ne reproduit pas le format officiel de l’EAM.
\begin{enumerate}
\item Les racines de $(x+1)(x-4)$ sont
\begin{enumerate}[label=\Alph*.,itemsep=2pt]
\item $1$ et $4$
\item $-1$ et $-4$
\item $-1$ et $4$
\item $1$ et $-4$
\end{enumerate}
\item Le discriminant de $x^2-2x+3$ vaut
\begin{enumerate}[label=\Alph*.,itemsep=2pt]
\item $-8$
\item $8$
\item $16$
\item $-2$
\end{enumerate}
\item Le sommet de $(x-2)^2+5$ est
\begin{enumerate}[label=\Alph*.,itemsep=2pt]
\item $\pt{2;-5}$
\item $\pt{-2;5}$
\item $\pt{5;2}$
\item $\pt{2;5}$
\end{enumerate}
\item La somme des racines de $2x^2-6x+1$ vaut
\begin{enumerate}[label=\Alph*.,itemsep=2pt]
\item $6$
\item $3$
\item $-3$
\item $\dfrac{1}{2}$
\end{enumerate}
\item $x^2-4<0$ équivaut à
\begin{enumerate}[label=\Alph*.,itemsep=2pt]
\item $x>2$
\item $x<-2$
\item $x\in\intoo{-2}{2}$
\item $x\in\intff{-2}{2}$
\end{enumerate}
\mxcounted{qcm1}
\end{enumerate}
\endgroup
\begingroup\small
\mxneed{18}\section*{QCM 2 : comprendre et choisir}\label{C02-Q2}
\mxCount{qcm2} questions pédagogiques à réponse unique. Choisir une réponse par question ; expliquer ensuite un distracteur. Ce format de chapitre ne reproduit pas le format officiel de l’EAM.
\begin{enumerate}
\item $x^2+6x+9$ a pour factorisation
\begin{enumerate}[label=\Alph*.,itemsep=2pt]
\item $(x-3)^2$
\item $(x+9)(x+1)$
\item $(x+6)^2$
\item $(x+3)^2$
\end{enumerate}
\item Le produit des racines de $3x^2+x-2$ vaut
\begin{enumerate}[label=\Alph*.,itemsep=2pt]
\item $-\dfrac{2}{3}$
\item $-\dfrac{1}{3}$
\item $\dfrac{2}{3}$
\item $-2$
\end{enumerate}
\item $-x^2+4x$ atteint sur $\intff{0}{1}$ son maximum en
\begin{enumerate}[label=\Alph*.,itemsep=2pt]
\item $2$
\item $0$
\item $1$
\item $4$
\end{enumerate}
\item Une racine double impose
\begin{enumerate}[label=\Alph*.,itemsep=2pt]
\item $a=0$
\item $\Delta=0$
\item $b=0$
\item $c=0$
\end{enumerate}
\item Les polynômes de degré 2 de racines 0 et 2 sont
\begin{enumerate}[label=\Alph*.,itemsep=2pt]
\item $a(x+2)$
\item $a(x-2)^2$
\item $x(x-2)$ uniquement
\item $ax(x-2)$, $a\ne0$
\end{enumerate}
\mxcounted{qcm2}
\end{enumerate}
\subsection*{Vrai ou faux ? Justifier.}\label{C02-VF}
\begin{enumerate}
\item Une fonction du second degré possède toujours deux racines réelles.
\item Un trinôme ayant deux racines distinctes est négatif entre elles.
\item Pour tout réel x, $(x-2)^2\geq0$.
\end{enumerate}
\begin{keypoint}[Bilan personnel]Je justifie le choix de la forme ; je distingue racines et sommet ; je termine une inéquation par un ensemble ; je respecte le domaine du modèle.

Pour chaque objectif, indiquer : acquis sans aide ; réussi avec aide ; à retravailler. Écrire un exemple et une erreur à éviter.\end{keypoint}
\endgroup
\ifmxteacher\mxinclude{corriges/C02}\fi
